\section{Discussion}
\label{sec:disc_maintenance}

Factors related to mobility, such as trip distance, not only differentiate electric from mechanical mobility (Chapter~\ref{ch:bicing_mobility_patterns}) but were also found to significantly influence the wear and tear on bike components. Leveraging these insights, we present a novel \ac{pm} system for a \ac{bss}, which marks a significant step towards replacing corrective maintenance strategies. The system has been designed to predict maintenance needs for three essential bike components, delivering satisfactory forecasting results through the application of \ac{dl} survival models. The design we introduce enables operation across an entire bike fleet, distinguishing it from previous research in bike \ac{pm} that primarily focused on individual bikes~\cite{MountainBike2019, Matkovic2021}. Moreover, our system enhances upon key aspects of an earlier \ac{pm} solution developed for Oslo's \ac{bss}~\cite{Predictivemaintenancebike2018}. First, our datasets are considerably more comprehensive. Second, we treat different repair typologies independently, acknowledging their distinct breakdown dynamics.  Lastly, to prevent potentially discriminatory conclusions that could impact BSS pricing strategies, we have deliberately omitted user information such as gender and age from our modeling. Additionally, through the application of a game-theoretic interpretability approach, we verify that the model's predictions are consistent with the observed failure dynamics. Notably, the most influential factors in the generation of the predictions are the cumulative distance and whether the bike is mechanical or electric.

The \ac{pm} results presented here, together with the mobility patterns documented in Chapter~\ref{ch:bicing_mobility_patterns}, hold significant promise for impacting society, especially in terms of enhancing the sustainability of urban mobility. Understanding \ac{bss} users' preferences for mechanical and electric mobility is essential for improving the decision-support systems that facilitate fleet rebalancing. Enhancing these systems not only assists \ac{bss} managers but also promotes \ac{bss} mobility by improving the user's experience, ensuring the availability of the preferred mode of transport when needed. On the other hand, the \ac{pm} system we presented, which can predict the breakdowns of three key bike components, marks an initial step towards developing a more comprehensive system that considers more bike components. Implementing such holistic systems could significantly enhance urban mobility sustainability by reducing the environmental footprint and operational costs of \ac{bss} through more efficient resource utilization. Beyond improving the scheduling of maintenance activities, these systems also facilitate the application of bike rebalancing strategies to extend the lifetime of bike parts. For instance, by strategically relocating bikes to less demanding routes or stations, we can minimize stress on critical components, thereby prolonging their usability when needed. Furthermore, the \ac{pm} systems would enhance the user's experience by ensuring a more reliable and available bike fleet. However, while this study demonstrates the feasibility of successfully deploying a \ac{pm} system for \ac{bss}, employing interpretability tools to build trust in the accuracy of these systems among \ac{bss} managers is crucial for their broader adoption. 

Although our data was of high quality, several limitations may have influenced our results. Because we lacked precise bike route information, we had to estimate trip distances, which may have impacted aspects of the \ac{pm} modeling—particularly for covariates such as cumulative distance. Similarly, we simplified elevation changes, potentially missing irregularities in the routes between stations that could affect wear-related variables. In addition, maintenance records indicated when components were replaced rather than when they actually failed. Finally, imbalances in the data due to censoring could have affected the predictive performance of our survival models.

Finally, building on the findings of this study, future research could extend the \ac{pm} framework to include more components, developing a more comprehensive system. A crucial step to validate this approach would involve the real-world deployment of these models, providing practical evidence of their impacts.
